\begin{lemma}[Block-to-letter translation of the boundary block matrix equation]
    \label{lem:block_mateq_of_blocked_mateq}
    \lean{MPSTensor.block_matEq_of_blocked_matEq}
    \leanok
    \uses{def:eval_word, def:blocked_tensor, lem:eval_word_block}
    Let $A$ be an MPS tensor and let $L_0,K\ge 0$. Let $X$ be a matrix, and let
    $Y_{c_b}$ be a matrix for each iterated block index $c_b$ of the complement.
    Suppose that, for every block letter $b$ of the alphabet
    $\Fin{d}^{L_0}$ and every iterated block index $c_b$ of the
    complement, with $w(b)$ the length-$L_0$ word of $b$ and
    $\widetilde{w(c_b)}$ the length-$L_0K$ complement word obtained by
    concatenating the blocks of $c_b$,
    $X A^{w(b)} A^{\widetilde{w(c_b)}} = A^{w(b)} Y_{c_b}$.
    Then there is a family $Y'_c$, indexed by the length-$L_0K$ words $c$, with
    $X A^s A^c = A^s Y'_c$
    for every length-$L_0$ word $s$ and every length-$L_0K$ word $c$.
\end{lemma}

\begin{proof}
    \leanok
    Each length-$L_0$ word $s$ is the word $w(b)$ of its block letter $b$,
    and each length-$L_0K$ word $c$ regroups into $K$ blocks, recovering it as
    the concatenated complement word $\widetilde{w(c_b)}$ of an iterated block
    index $c_b$; both identifications are bijective. Substituting these
    representations into the hypothesis and setting $Y'_c=Y_{c_b}$ gives the
    asserted equation for every $s$ and $c$.
\end{proof}

\begin{lemma}[Boundary matrix commutation from block-window equations]
    \label{lem:boundary_matrix_commutes_block_window}
    \lean{MPSTensor.boundary_matrix_commutes_of_isNBlkInjective_of_block_matEq}
    \leanok
    \uses{def:l_blk_injective, def:eval_word, lem:wordspan_blk_inj,
        lem:padded_complement_annihilation, lem:block_mateq_of_blocked_mateq}
    Let $A$ be $L_0$-block-injective with $L_0>0$. Let $X$ be a boundary
    matrix, and let $Y_v$ be a matrix for each complementary word $v$ of
    length $K$. Suppose that, for every word $u$ of length $L_0$ and every
    word $v$ of length $K$,
    $X A^u A^v = A^u Y_v$.
    Then, for every physical letter $j$,
    $X A^j = A^j X$.
    This is the boundary-matrix commutation step obtained from the equations
    $X A^u A^v=A^uY_v$, for all length-$L_0$ words $u$ and all length-$K$
    words $v$, as in
    \cite[Section~IV.C, lines~2049--2090]{Cirac2021Matrix}.
\end{lemma}

\begin{proof}
    \leanok
    \uses{lem:wordspan_blk_inj, lem:padded_complement_annihilation}
    By Lemma~\ref{lem:wordspan_blk_inj},
    $\spn\{A^u : |u|=L_0\} = \MN{D}$.
    Hence, for every $M\in\MN{D}$ and every complementary word $v$ of length
    $K$, the displayed hypothesis extends by linearity in the length-$L_0$
    block to $XMA^v=MY_v$. Taking $M=\Id$ gives
    \begin{align}
        Y_v &= X A^v.
        \label{eq:ph_normal_closing_word_value}
    \end{align}
    Taking $M=A^j$ and substituting (\ref{eq:ph_normal_closing_word_value})
    gives $(X A^j-A^jX)A^v=0$ for every word $v$ of length $K$.
    Since $K\le (K+1)L_0$ and $K+1\ge1$,
    Lemma~\ref{lem:padded_complement_annihilation} gives
    $X A^j-A^jX=0$.
\end{proof}

\begin{lemma}[Trace identity from the boundary-crossing cyclic window]
    \label{lem:closure_property_boundary_block_window_trace_ground_space_map}
    \lean{MPSTensor.closure_property_boundary_block_window_trace_eq_of_groundSpaceMap}
    \leanok
    \uses{def:ground_space_map, rem:cyclic_window_operators}
    Let $A$ be a tensor with virtual dimension at least one, and let
    $L_0>0$ and $L_0<M$. Let $\psi=\Gamma_{M+1}(X)$ and assume that every
    cyclic restriction of length $L_0+1$ belongs to $G_{L_0+1}(A)$. Then there
    are boundary matrices $Y_\nu$, indexed by nonempty complementary words
    $\nu$ of length $M+1-(L_0+1)$, such that, for every physical letter $j$
    and every word $\alpha$ of length $L_0$,
    \begin{align}
        \tr\left(A^j X A^\alpha A^\nu\right)
        &= \tr\left(A^j A^\alpha Y_\nu\right).
        \notag
    \end{align}
\end{lemma}

\begin{proof}
    \leanok
    Let $Y_\nu$ be the matrix representing the cyclic restriction whose window
    begins at the last site. The cyclic word is $\alpha\,\nu\,j$ after deleting
    the first letter from the window and reading the complement, so the
    ground-space representation and cyclicity of the trace give
    \begin{align}
        \tr\left(A^jX A^\alpha A^\nu\right)
        &= \tr\left(A^\alpha A^\nu A^jX\right)
         = \tr\left(A^jA^\alpha Y_\nu\right).
        \notag
    \end{align}
\end{proof}

\begin{lemma}[One-site-injective boundary block-window identity]
    \label{lem:closure_property_boundary_block_window_equation_ground_space_map_injective}
    \lean{MPSTensor.closure_property_boundary_block_window_equation_of_groundSpaceMap_of_isInjective}
    \leanok
    \uses{def:injective, rem:cyclic_window_operators,
        lem:closure_property_boundary_block_window_trace_ground_space_map}
    Let $A$ be a tensor with virtual dimension at least one whose single-site
    matrices span the full matrix algebra. Let $L_0>0$ and $L_0<M$, and let
    $\psi=\Gamma_{M+1}(X)$. If every cyclic restriction of length $L_0+1$
    belongs to $G_{L_0+1}(A)$, then there are boundary matrices $Y_\nu$,
    indexed by nonempty complementary words $\nu$ of length
    $M+1-(L_0+1)$, such that, for every word $\alpha$ of length $L_0$,
    $X A^\alpha A^\nu = A^\alpha Y_\nu$.
    % Scope restriction: the source argument in arXiv:2011.12127, Section~IV.C,
    % lines 2078--2079, assumes only $L_0$-block injectivity. This lemma records
    % the narrower one-site-injective case; the general case is documented in
    % docs/paper-gaps/cpgsv21_normal_range_reduction.tex.
\end{lemma}

\begin{proof}
    \leanok
    \uses{lem:closure_property_boundary_block_window_trace_ground_space_map,
        thm:trace_pairing_injective}
    For every physical letter $j$,
    Lemma~\ref{lem:closure_property_boundary_block_window_trace_ground_space_map}
    gives
    \begin{align}
        \tr\left(A^jX A^\alpha A^\nu\right)
        &= \tr\left(A^jA^\alpha Y_\nu\right).
        \label{eq:ph_normal_closing_trace_window}
    \end{align}
    Put $Z=X A^\alpha A^\nu-A^\alpha Y_\nu$. For every physical letter $j$,
    Then (\ref{eq:ph_normal_closing_trace_window}) gives $\tr(A^jZ)=0$.
    Since the matrices $A^j$ span the full matrix algebra, the trace pairing
    gives
    $(\forall j,\ \tr(A^jZ)=0) \quad\Longrightarrow\quad Z=0$.
    Hence $X A^\alpha A^\nu=A^\alpha Y_\nu$.
\end{proof}

\begin{theorem}[Trace identities with length-$L_0$ probes from cyclic-window constraints]
    \label{thm:closure_property_boundary_block_window_trace_word_ground_space_map}
    \lean{MPSTensor.closure_property_boundary_block_window_trace_evalWord_mul_eq_of_groundSpaceMap}
    \leanok
    \uses{def:ground_space_map, rem:cyclic_window_operators,
        lem:closure_property_boundary_crossing_equations_ground_space_map,
        lem:closure_property_boundary_condition_product_window_witnesses,
        lem:outside_labels_determine_boundary_witnesses}
    Let $A$ be $L_0$-block-injective with $L_0>0$, $L_0<M$, and
    $D\ge1$. Let $\psi=\Gamma_{M+1}(X)$ and assume that every cyclic
    restriction of length $L_0+1$ belongs to $G_{L_0+1}(A)$. Then there are
    boundary matrices $Y_\nu$, indexed by nonempty complementary words $\nu$
    of length $M+1-(L_0+1)$, such that, for every word $\alpha$ of length
    $L_0$ and every word $\beta$ of length $L_0$,
    \begin{align}
        \tr\left(A^\beta X A^\alpha A^\nu\right)
        &= \tr\left(A^\beta A^\alpha Y_\nu\right).
        \notag
    \end{align}
\end{theorem}

\begin{proof}
    \leanok
    Choose representation matrices for every cyclic restriction of length
    $L_0+1$. Fix $\alpha$ and $\nu$. The boundary-crossing equations move the
    matrix $X$ through the first boundary letter. In local boundary notation,
    this has the form
    \begin{align}
        X A^{\rho_0}A^{\rho_1\cdots\rho_{M-L_0}}
        &= A^{\rho_0}Y_{M+1-L_0}(\rho).
        \notag
    \end{align}
    The adjacent boundary-window product then transports this comparison
    through the remaining $L_0-1$ boundary letters:
    \begin{align}
        Y_{M+1-L_0}(\rho)A^{\rho_{M+1-L_0}\cdots\rho_{M-1}}
        &= A^{\rho_1\cdots\rho_{L_0-1}}Y_M(\rho).
        \notag
    \end{align}
    Finally, the outside-label uniqueness lemma identifies the final boundary
    matrix with the matrix $Y_\nu$ attached to the complementary word $\nu$.
    Thus $X A^\alpha A^\nu=A^\alpha Y_\nu$. Multiplying on the left by
    $A^\beta$ and taking the trace gives the asserted identity for every
    length-$L_0$ word $\beta$.
\end{proof}

\begin{theorem}[Length-word trace separation for the boundary block window]
    \label{thm:block_window_matrix_equation_trace_word_products}
    \lean{MPSTensor.block_window_matrix_equation_of_trace_evalWord_mul_eq_of_isNBlkInjective}
    \leanok
    \uses{def:l_blk_injective, def:eval_word}
    Let $A$ be $L_0$-block-injective. Let $X$ be a boundary matrix, and let
    $Y_\nu$ be a matrix for each complementary word $\nu$ of length $K$.
    Suppose that, for every word $\alpha$ of length $L_0$, every word $\beta$
    of length $L_0$, and every complementary word $\nu$,
    \begin{align}
        \tr\left(A^\beta X A^\alpha A^\nu\right)
        &= \tr\left(A^\beta A^\alpha Y_\nu\right).
        \notag
    \end{align}
    Then, for every word $\alpha$ of length $L_0$ and every complementary word
    $\nu$,
    $X A^\alpha A^\nu = A^\alpha Y_\nu$.
\end{theorem}

\begin{proof}
    \leanok
    \uses{thm:ground_space_map_injective_blk}
    The trace identities say exactly that, for each fixed pair
    $(\alpha,\nu)$,
    \begin{align}
        \Gamma_{L_0}(X A^\alpha A^\nu)
        &= \Gamma_{L_0}(A^\alpha Y_\nu).
        \notag
    \end{align}
    Since $A$ is $L_0$-block-injective, $\Gamma_{L_0}$ is injective by
    Theorem~\ref{thm:ground_space_map_injective_blk}. Hence
    $X A^\alpha A^\nu=A^\alpha Y_\nu$.
\end{proof}

\begin{lemma}[Matrix identity $X A^\alpha A^\nu=A^\alpha Y_\nu$ from cyclic-window constraints]
    \label{lem:closure_property_boundary_block_window_equation_ground_space_map}
    \lean{MPSTensor.closure_property_boundary_block_window_equation_of_groundSpaceMap}
    \leanok
    \uses{def:ground_space_map, rem:cyclic_window_operators,
        thm:closure_property_boundary_block_window_trace_word_ground_space_map,
        thm:block_window_matrix_equation_trace_word_products}
    Let $A$ be $L_0$-block-injective with $L_0>0$, $L_0<M$, and
    $D\ge1$. Let $\psi=\Gamma_{M+1}(X)$ and assume that every cyclic
    restriction of length $L_0+1$ belongs to $G_{L_0+1}(A)$. Then there are
    boundary matrices $Y_\nu$, indexed by nonempty complementary words $\nu$
    of length $M+1-(L_0+1)$, such that, for every word $\alpha$ of length
    $L_0$,
    $X A^\alpha A^\nu = A^\alpha Y_\nu$.
\end{lemma}

\begin{proof}
    \leanok
    Theorem~\ref{thm:closure_property_boundary_block_window_trace_word_ground_space_map}
    gives, for every word $\beta$ of length $L_0$,
    \begin{align}
        \tr\left(A^\beta X A^\alpha A^\nu\right)
        &= \tr\left(A^\beta A^\alpha Y_\nu\right).
        \notag
    \end{align}
    Theorem~\ref{thm:block_window_matrix_equation_trace_word_products} then
    gives the displayed matrix equation.
\end{proof}

\begin{lemma}[Boundary restrictions from commutation and one-sided products]
    \label{lem:boundary_restrictions_from_commutation_one_sided_products}
    \lean{MPSTensor.boundary_restrictions_eq_of_commutes_and_one_sided}
    \leanok
    \uses{def:l_blk_injective, def:eval_word,
        lem:one_sided_boundary_witness_unique}
    Let $A$ be $L_0$-block-injective with $L_0>0$. Let $X$ be a boundary
    matrix, let $W_\eta$ and $V_\eta$ be matrix families indexed by a physical
    letter $\eta$, and let $\mu$ be a word. Suppose that, for every physical
    letter $j$,
    $X A^j = A^jX$.
    Suppose also that, for every $\eta$ and every $j$,
    \begin{align}
        W_\eta A^j &= A^\mu A^jX,
        \notag\\
        X A^jA^\mu &= A^jV_\eta.
        \notag
    \end{align}
    Then, for every $\eta$ and every $j$,
    $W_\eta A^j = V_\eta A^j$.
\end{lemma}

\begin{proof}
    \leanok
    \uses{lem:one_sided_boundary_witness_unique}
    The commutation relation $X A^j=A^jX$ extends by induction to
    $X A^w=A^wX$ for every word $w$. Hence the second one-sided equation gives,
    for every $k$,
    $A^kV_\eta = X A^kA^\mu = A^kX A^\mu = A^kA^\mu X$.
    Lemma~\ref{lem:one_sided_boundary_witness_unique} gives
    $V_\eta=A^\mu X$. The first one-sided equation then gives
    $W_\eta A^j = A^\mu A^jX = A^\mu X A^j = V_\eta A^j$.
\end{proof}

\begin{lemma}[Products with one-site tensors determine the boundary matrix]
    \label{lem:wrapped_mirror_right_products_determine}
    \lean{MPSTensor.wrapped_mirror_witness_agree_of_right_products}
    \leanok
    \uses{lem:one_sided_boundary_witness_unique}
    Let $A$ be $L_0$-block-injective with $L_0>0$. Fix two matrix
    families $Y^+$ and $Y^-$ and the two reindexed boundary conditions
    $\tau^+_\eta(\mu)$ and $\tau^-_\eta(\mu)$. If, for every physical letter
    $j$,
    $Y^+_{\tau^+_\eta(\mu)} A^j = Y^-_{\tau^-_\eta(\mu)} A^j$,
    then
    $Y^+_{\tau^+_\eta(\mu)} = Y^-_{\tau^-_\eta(\mu)}$.
\end{lemma}

\begin{proof}
    \leanok
    Apply Lemma~\ref{lem:one_sided_boundary_witness_unique} to the two
    matrices $Y^+_{\tau^+_\eta(\mu)}$ and $Y^-_{\tau^-_\eta(\mu)}$.
\end{proof}

\begin{lemma}[Outside labels determine boundary-restriction matrices]
    \label{lem:outside_labels_determine_boundary_witnesses}
    \lean{MPSTensor.wrappedMiddleBackground_witness_eq_of_complement_eq}
    \lean{MPSTensor.mirrorMiddleBackground_witness_eq_of_complement_eq}
    \leanok
    \uses{lem:one_sided_boundary_witness_unique, rem:cyclic_window_operators}
    Let $A$ be $L_0$-block-injective with $L_0>0$ and $L_0\le M$. If an outside
    configuration $\rho$ has the same word $\mu$ as
    $\tau^+_\eta(\mu)$ on the sites outside the last-site cyclic window, then
    the two matrices representing that restriction agree:
    $Y_\rho = Y_{\tau^+_\eta(\mu)}$.
    The same assertion holds for the second boundary-crossing window: if
    $\rho$ has the same outside word as $\tau^-_\eta(\mu)$, then
    $Y_\rho = Y_{\tau^-_\eta(\mu)}$.
    This is the boundary-matrix independence step used to choose one
    representative outside configuration for each outside word in the
    closure-property matrix comparison.
\end{lemma}

\begin{proof}
    \leanok
    Equality of the cyclic restrictions gives, for every physical letter $j$,
    $Y_\rho A^j = Y_{\tau^\pm_\eta(\mu)}A^j$.
    Lemma~\ref{lem:one_sided_boundary_witness_unique} identifies the two
    matrices from these one-site products.
\end{proof}

\begin{lemma}[First-letter restrictions determine a vector]
    \label{lem:first_letter_restrictions_determine}
    \lean{MPSTensor.eq_of_forall_restrictFirst_eq}
    \leanok
    Let $R_j$ be restriction to first letter $j$:
    \begin{align}
        (R_j\phi)(\sigma_1,\ldots,\sigma_L)
        &= \phi(j,\sigma_1,\ldots,\sigma_L).
        \notag
    \end{align}
    For $\phi,\psi\in(\C^d)^{\otimes(L+1)}$,
    \begin{align}
        (\forall j,\ R_j\phi=R_j\psi)
        &\quad\Longrightarrow\quad \phi=\psi.
        \notag
    \end{align}
\end{lemma}

\begin{proof}
    \leanok
    At $(\sigma_0,\sigma_1,\ldots,\sigma_L)$, the required equality is
    \begin{align}
        \phi(\sigma_0,\sigma_1,\ldots,\sigma_L)
        &= (R_{\sigma_0}\phi)(\sigma_1,\ldots,\sigma_L)
         = (R_{\sigma_0}\psi)(\sigma_1,\ldots,\sigma_L)
         = \psi(\sigma_0,\sigma_1,\ldots,\sigma_L).
        \notag
    \end{align}
\end{proof}

\begin{lemma}[Auxiliary restriction equality for the closure property]
    \label{lem:boundary_restriction_from_fixed_letters}
    \lean{MPSTensor.closure_property_boundary_restriction_eq_of_fixed_boundary_letters}
    \leanok
    \uses{lem:first_letter_restrictions_determine, rem:cyclic_window_operators}
    Let $L_0>0$ and $L_0\le M$, and set
    \begin{align}
        B^+_{\eta,\mu}(\psi)
        &= \Res^{\tau^+_{\eta}(\mu)}_{M,L_0+1}(\psi),
        \notag\\
        B^-_{\eta,\mu}(\psi)
        &= \Res^{\tau^-_{\eta}(\mu)}_{M+1-L_0,L_0+1}(\psi).
        \notag
    \end{align}
    With $R_j$ as in Lemma~\ref{lem:first_letter_restrictions_determine},
    \begin{align}
        \bigl(\forall j,\ R_jB^+_{\eta,\mu}(\psi)
        = R_jB^-_{\eta,\mu}(\psi)\bigr)
        &\quad\Longrightarrow\quad
        B^+_{\eta,\mu}(\psi)=B^-_{\eta,\mu}(\psi).
        \notag
    \end{align}
\end{lemma}

\begin{proof}
    \leanok
    Apply Lemma~\ref{lem:first_letter_restrictions_determine} to the two
    length-$(L_0+1)$ restrictions, using the displayed family of first-letter
    equalities.
\end{proof}

\begin{lemma}[Boundary-crossing restriction from right products]
    \label{lem:closed_boundary_restriction_from_right_products}
    \lean{MPSTensor.closure_property_boundary_restriction_eq_of_first_products}
    \leanok
    \uses{rem:cyclic_window_operators, lem:first_letter_restrictions_determine}
    Let $A$ be an MPS tensor, let $L_0>0$, and let $L_0\le M$. Suppose the two
    boundary-crossing restrictions of length $L_0+1$ are represented by boundary
    matrices $Y^+_{\eta,\mu}$ and $Y^-_{\eta,\mu}$:
    \begin{align}
        B^+_{\eta,\mu}(\psi)
        &= \Gamma_{L_0+1}(Y^+_{\eta,\mu}),
        \notag\\
        B^-_{\eta,\mu}(\psi)
        &= \Gamma_{L_0+1}(Y^-_{\eta,\mu}).
        \notag
    \end{align}
    If, for every physical letter $j$,
    $Y^+_{\eta,\mu}A^j = Y^-_{\eta,\mu}A^j$,
    then
    $B^+_{\eta,\mu}(\psi) = B^-_{\eta,\mu}(\psi)$.
\end{lemma}

\begin{proof}
    \leanok
    Fix $j$ and take the first-letter restriction of the two displayed
    length-$(L_0+1)$ restrictions. These are
    $\Gamma_{L_0}(Y^+_{\eta,\mu}A^j)$ and
    $\Gamma_{L_0}(Y^-_{\eta,\mu}A^j)$. The one-site product equality gives,
    for every $j$,
    \begin{align}
        \Gamma_{L_0}(Y^+_{\eta,\mu}A^j)
        &= \Gamma_{L_0}(Y^-_{\eta,\mu}A^j).
        \notag
    \end{align}
    Lemma~\ref{lem:first_letter_restrictions_determine} identifies the two
    original restrictions.
\end{proof}

\begin{lemma}[Boundary restrictions from first-letter products]
    \label{lem:closed_boundary_restrictions_from_first_letter_products}
    \lean{MPSTensor.closure_property_boundary_restrictions_eq_of_first_products}
    \leanok
    \uses{lem:closed_boundary_restriction_from_right_products}
    Let $L_0>0$ and $L_0\le M$. Let $Y_i(\tau)$ represent the
    length-$(L_0+1)$ cyclic restrictions of $\psi$. Fix a complementary
    word $\mu$. If, for every boundary letter $\eta$ and physical letter
    $j$,
    $Y_M(\tau^+_\eta(\mu))A^j = Y_{M+1-L_0}(\tau^-_\eta(\mu))A^j$,
    then, for every $\eta$,
    \begin{align}
        \Res^{\tau^+_\eta(\mu)}_{M,L_0+1}(\psi)
        &= \Res^{\tau^-_\eta(\mu)}_{M+1-L_0,L_0+1}(\psi).
        \notag
    \end{align}
\end{lemma}

\begin{proof}
    \leanok
    Apply Lemma~\ref{lem:closed_boundary_restriction_from_right_products} for
    each boundary letter $\eta$, using the corresponding first-letter product
    equations.
\end{proof}

\begin{lemma}[Equality of the two restrictions crossing the periodic cut]
    \label{lem:closure_property_boundary_restriction_chain_ground_space}
    \lean{MPSTensor.closure_property_boundary_restriction_eq_of_chainGroundSpace}
    \leanok
    \uses{def:chain_ground_space, rem:cyclic_window_operators,
        lem:chain_ground_space_window_witnesses,
        lem:closed_boundary_restriction_from_right_products,
        thm:boundary_closing_right_products_chain_ground_space}
    Let $A$ be $L_0$-block-injective with $L_0>0$ and $D\ge1$, let
    $L_0<M$, and let
    \begin{align}
        \psi &\in\mathcal G_{M+1,L_0+1}(A),
        \notag\\
        \psi &= \Gamma_{M+1}(X).
        \notag
    \end{align}
    For every boundary letter $\eta$ and complementary word $\mu$, the two
    restrictions crossing the periodic cut satisfy
    \begin{align}
        \Res^{\tau^+_{\eta}(\mu)}_{M,L_0+1}(\psi)
        &= \Res^{\tau^-_{\eta}(\mu)}_{M+1-L_0,L_0+1}(\psi).
        \notag
    \end{align}
\end{lemma}

\begin{proof}
    \leanok
    The cyclic-window representation lemma gives matrices
    $Y_M(\tau^+_\eta(\mu))$ and
    $Y_{M+1-L_0}(\tau^-_\eta(\mu))$ such that
    \begin{align}
        \Res^{\tau^+_\eta(\mu)}_{M,L_0+1}(\psi)
        &= \Gamma_{L_0+1}\!\left(Y_M(\tau^+_\eta(\mu))\right),
        \notag\\
        \Res^{\tau^-_\eta(\mu)}_{M+1-L_0,L_0+1}(\psi)
        &= \Gamma_{L_0+1}\!\left(Y_{M+1-L_0}(\tau^-_\eta(\mu))\right).
        \notag
    \end{align}
    Theorem~\ref{thm:boundary_closing_right_products_chain_ground_space}
    gives, for every $j$,
    $Y_M(\tau^+_\eta(\mu))A^j = Y_{M+1-L_0}(\tau^-_\eta(\mu))A^j$.
    Lemma~\ref{lem:closed_boundary_restriction_from_right_products} identifies
    the two length-$(L_0+1)$ restrictions.
\end{proof}
